%========================================
% MatinBook — Stage 15: Cover and Full Book Test
% Version: v1.1
% Purpose:
%   Validate the cover system and the complete book assembly:
%     1. Front cover (\makecover)
%     2. Back cover (\makebackcover + \authorbio)
%     3. Full book structure (frontmatter, mainmatter, backmatter)
%     4. Bibliography (biblatex + biber)
%     5. Index (xindy with Persian sorting)
%     6. Cover and back-cover interaction with the rest of
%        the book
%
% Compilation:
%   xelatex -shell-escape stage15-cover.tex
%   biber stage15-cover
%   xelatex -shell-escape stage15-cover.tex
%   xelatex -shell-escape stage15-cover.tex
%   (full cycle required for bibliography + index + cover)
%
% Expected outcome:
%   A complete book with a professionally typeset front
%   cover, a well-structured back cover with book blurb and
%   author bio, and all the internal machinery (ToC, lists,
%   bibliography, index) working correctly.
%========================================

%------------------------------------------------------------------------
% Search Path (for tests directory)
%------------------------------------------------------------------------
\makeatletter
\def\input@path{{../}}
\makeatother

%------------------------------------------------------------------------
% Document Class
%------------------------------------------------------------------------
\documentclass{matinbook}

%------------------------------------------------------------------------
% Bibliography Resource
%------------------------------------------------------------------------
\addbibresource{../references.bib}

%------------------------------------------------------------------------
% Book Metadata
%------------------------------------------------------------------------
\title{هنر برنامه‌نویسی}
\author{تیم توسعه MatinBook}
\date{مهر ۱۴۰۵}

\booktitle{هنر برنامه‌نویسی}

\renewcommand{\repository}{github.com/matinmahmoudi-cs/matinbook}

%========================================================================
% DOCUMENT
%========================================================================
\begin{document}

%------------------------------------------------------------------------
% FRONT COVER
%------------------------------------------------------------------------
\makecover
    {هنر برنامه‌نویسی}
    {از ایده تا اجرا}
    {تیم توسعه MatinBook}
    {مهر ۱۴۰۵}

%------------------------------------------------------------------------
% FRONT MATTER (symmetric margins)
%------------------------------------------------------------------------
\frontmatter
\frontmattergeometry

% Title page
\maketitle

% Copyright / Colophon
\thispagestyle{empty}
\vfill
\begin{center}
    {\large\textbf{حقوق نشر}}

    \vspace{1cm}

    تمامی حقوق این کتاب برای تیم \lr{MatinBook} محفوظ است.
    هرگونه کپی‌برداری، تکثیر و انتشار بدون اجازه‌ی کتبی
    ممنوع است.

    \vspace{1cm}

    نسخه \matinbookversion\ — \matinbookdate

    \vspace{1cm}

    \lr{Copyright \copyright\ 2026 MatinBook Project} \\
    \lr{Licensed under the MIT License}
\end{center}
\clearpage

% Dedication
\thispagestyle{empty}
\vfill
\begin{flushright}
    {\Large\itshape
        به تمام کسانی که\\
        با شوق می‌آموزند\\
        و با عشق می‌آموزانند.
    }
\end{flushright}
\clearpage

% Preface
\chapter*{پیشگفتار}
\addcontentsline{toc}{chapter}{پیشگفتار}

این کتاب برای کسانی نوشته شده که می‌خواهند
برنامه‌نویسی را به‌صورت اصولی بیاموزند. مطالب از
سطح پایه شروع شده و به‌تدریج به مباحث پیشرفته
می‌رسد.

امیدواریم این کتاب بتواند سهمی هرچند کوچک در ارتقای
دانش برنامه‌نویسی فارسی‌زبانان داشته باشد.

\vspace{1cm}
\begin{flushleft}
    \textbf{تیم توسعه MatinBook} \\
    مهر ۱۴۰۵
\end{flushleft}
\clearpage

% Table of Contents
\tableofcontents
\clearpage

% List of Figures
\listoffigures
\clearpage

% List of Tables
\listoftables
\clearpage

% List of Algorithms
\listofalgorithms
\clearpage

%------------------------------------------------------------------------
% MAIN MATTER (wide margin for notes)
%------------------------------------------------------------------------
\mainmatter
\mainmattergeometry

%========================================================================
% CHAPTER 1 — INTRODUCTION
%========================================================================
\chapter{مقدمه}
\label{chap:introduction}

\section{چرا برنامه‌نویسی؟}
\label{sec:intro-why}

برنامه‌نویسی\index{برنامه‌نویسی} مهارتی است که در
دنیای امروز به اندازه‌ی خواندن و نوشتن اهمیت دارد.
از تلفن‌های همراه گرفته تا خودروهای خودران، همه با
کد کار می‌کنند.

\mnote{برنامه‌نویسی یکی از مهارت‌های کلیدی قرن بیست‌و‌یکم است.}

\begin{definition}[برنامه‌نویسی]
    برنامه‌نویسی فرایند طراحی، نوشتن، آزمایش و
    نگهداری کدهای کامپیوتری است که به ماشین می‌گویند
    چه کاری انجام دهد.
    \label{def:programming}
\end{definition}

\section{انتخاب زبان برنامه‌نویسی}
\label{sec:intro-language}

برای شروع، پایتون\index{پایتون} بهترین انتخاب است.
دلایل:

\begin{itemize}
    \item نحو ساده و خوانا.
    \item جامعه‌ی کاربری بزرگ.
    \item کتابخانه‌های فراوان.
    \item کاربرد در حوزه‌های مختلف.
\end{itemize}

\begin{example}[اولین برنامه]
    ساده‌ترین برنامه در پایتون:

    \begin{latin}
\begin{minted}{python}
print("Hello, World!")
\end{minted}
    \end{latin}
    \label{ex:hello-world}
\end{example}

%========================================================================
% CHAPTER 2 — BASICS
%========================================================================
\chapter{مفاهیم پایه}
\label{chap:basics}

\section{متغیرها}
\label{sec:basics-variables}

متغیر\index{متغیر} محلی برای ذخیره‌ی داده است. در
پایتون، متغیرها نیاز به اعلان نوع ندارند
(\cref{code:variables}):

\begin{latin}
\begin{minted}{python}
# Variables in Python
name = "Ali"           # String
age = 25               # Integer
pi = 3.14              # Float
is_active = True       # Boolean

print(f"{name} is {age} years old")
\end{minted}
\end{latin}
\codecaption{تعریف متغیر در پایتون}
\label{code:variables}

\section{عملگرها}
\label{sec:basics-operators}

عملگرها\index{عملگر} عملیات روی داده‌ها را انجام
می‌دهند. جدول \cref{tab:operators} عملگرهای پایه را
نشان می‌دهد:

\begin{matintable}{عملگرهای پایه در پایتون}{tab:operators}
\begin{tblr}{
    width = \textwidth,
    colspec = {l X[c] X[c]},
    hlines, vlines,
    colsep = 8pt,
    rowsep = 4pt,
    row{1} = {font=\bfseries},
}
نوع & عملگر & مثال \\
محاسباتی & \lr{+, -, *, /, //, \%, **} & \lr{5 + 3 = 8} \\
مقایسه‌ای & \lr{==, !=, <, >, <=, >=} & \lr{5 > 3 = True} \\
منطقی & \lr{and, or, not} & \lr{True and False = False} \\
\end{tblr}
\end{matintable}

%========================================================================
% CHAPTER 3 — CONTROL FLOW
%========================================================================
\chapter{کنترل جریان}
\label{chap:controlflow}

\section{دستورات شرطی}
\label{sec:controlflow-conditional}

نمونه‌ای از دستور شرطی در \cref{code:if-else}:

\begin{latin}
\begin{minted}{python}
score = 85

if score >= 90:
    grade = "A"
elif score >= 80:
    grade = "B"
elif score >= 70:
    grade = "C"
else:
    grade = "F"

print(f"Score: {score}, Grade: {grade}")
\end{minted}
\end{latin}
\codecaption{دستور شرطی در پایتون}
\label{code:if-else}

\section{حلقه‌ها}
\label{sec:controlflow-loops}

حلقه‌های \lr{for} و \lr{while} در
\cref{code:loops}:

\begin{latin}
\begin{minted}{python}
# For loop
for i in range(5):
    print(f"Iteration {i}")

# While loop
count = 0
while count < 5:
    print(f"Count: {count}")
    count += 1
\end{minted}
\end{latin}
\codecaption{حلقه‌های \lr{for} و \lr{while}}
\label{code:loops}

%========================================================================
% CHAPTER 4 — FUNCTIONS
%========================================================================
\chapter{توابع}
\label{chap:functions}

\section{تعریف و استفاده}
\label{sec:functions-definition}

تابع\index{تابع} بلوکی از کد است که یک وظیفه‌ی مشخص
را انجام می‌دهد (\cref{code:functions}).

\begin{latin}
\begin{minted}{python}
def greet(name: str) -> str:
    """Return a greeting message."""
    return f"Hello, {name}!"

def add(a: int, b: int) -> int:
    """Return the sum of two numbers."""
    return a + b

# Using functions
print(greet("Ali"))
print(f"3 + 4 = {add(3, 4)}")
\end{minted}
\end{latin}
\codecaption{تعریف و استفاده از توابع}
\label{code:functions}

\begin{exercise}
    تابعی بنویسید که یک عدد دریافت کند و فاکتوریل آن
    را محاسبه کند.
    \label{exc:factorial}
\end{exercise}

\begin{solution}
    \begin{latin}
\begin{minted}{python}
def factorial(n: int) -> int:
    """Calculate factorial recursively."""
    if n <= 1:
        return 1
    return n * factorial(n - 1)

print(factorial(5))  # 120
\end{minted}
    \end{latin}
\end{solution}

%========================================================================
% CHAPTER 5 — DATA STRUCTURES
%========================================================================
\chapter{ساختمان داده‌ها}
\label{chap:datastructures}

\section{لیست}
\label{sec:datastructures-list}

لیست\index{لیست} پرکاربردترین ساختمان داده در
پایتون است (\cref{code:lists}).

\begin{latin}
\begin{minted}{python}
# Create and manipulate lists
numbers = [1, 2, 3, 4, 5]
numbers.append(6)        # Add to end
numbers.insert(0, 0)     # Insert at index
numbers.remove(3)        # Remove by value

print(f"List: {numbers}")
print(f"Length: {len(numbers)}")
print(f"Sum: {sum(numbers)}")
\end{minted}
\end{latin}
\codecaption{کار با لیست در پایتون}
\label{code:lists}

\section{دیکشنری}
\label{sec:datastructures-dict}

دیکشنری\index{دیکشنری} داده‌ها را به‌صورت
کلید-مقدار\index{کلید-مقدار} ذخیره می‌کند
(\cref{code:dictionary}).

\begin{latin}
\begin{minted}{python}
# Create a dictionary
student = {
    "name": "Ali",
    "age": 25,
    "grades": [18, 17, 20],
    "is_active": True
}

# Access and modify
print(student["name"])
student["age"] = 26
student["major"] = "Computer Science"

for key, value in student.items():
    print(f"{key}: {value}")
\end{minted}
\end{latin}
\codecaption{کار با دیکشنری}
\label{code:dictionary}

%========================================================================
% CHAPTER 6 — ALGORITHMS
%========================================================================
\chapter{الگوریتم‌ها}
\label{chap:algorithms}

\section{جستجوی دودویی}
\label{sec:algorithms-binary}

جستجوی دودویی (\cref{alg:binary}):

\begin{latin}
\begin{algorithm}
\caption{جستجوی دودویی}
\label{alg:binary}
\begin{algorithmic}[1]
    \Require Sorted array $A$, target $x$
    \Ensure Index of $x$ or $-1$
    \State $left \gets 0$, $right \gets n-1$
    \While{$left \leq right$}
        \State $mid \gets (left + right) // 2$
        \If{$A[mid] = x$}
            \State \Return $mid$
        \ElsIf{$A[mid] < x$}
            \State $left \gets mid + 1$
        \Else
            \State $right \gets mid - 1$
        \EndIf
    \EndWhile
    \State \Return $-1$
\end{algorithmic}
\end{algorithm}
\end{latin}

\begin{theorem}
    جستجوی دودویی (\cref{alg:binary}) در بدترین حالت
    \lr{$O(\log n)$} مقایسه انجام می‌دهد.
    \label{thm:binary}
\end{theorem}

%========================================================================
% CHAPTER 7 — VISUALIZATION
%========================================================================
\chapter{مصورسازی}
\label{chap:visualization}

\section{نمودار تابع}
\label{sec:visualization-plot}

نمودار \cref{fig:functions} دو تابع پایه را مقایسه
می‌کند:

\begin{figure}[h]
\centering
\begin{latin}
\begin{tikzpicture}
  \begin{axis}[
    width=\textwidth, height=6cm,
    xlabel={$x$},
    ylabel={$f(x)$},
    title={Common Functions},
    grid=major,
    legend pos=north west,
    domain=-2:2,
    samples=100
  ]
    \addplot[matinblue, thick] {x^2};
    \addlegendentry{$x^2$}
    \addplot[matinred, thick] {x^3};
    \addlegendentry{$x^3$}
  \end{axis}
\end{tikzpicture}
\end{latin}
\caption{نمودار $x^{2}$ و $x^{3}$}
\label{fig:functions}
\end{figure}

\section{فلوچارت}
\label{sec:visualization-flowchart}

فلوچارت \cref{fig:flowchart} یک فرآیند ساده را
نشان می‌دهد:

\begin{figure}[h]
\centering
\begin{latin}
\begin{tikzpicture}[
  node distance=2cm,
  scale=0.95,
  transform shape,
  arrow/.style={thick,->,>=stealth}
]
  \node[block] (start) {Start};
  \node[block, below of=start] (input) {Read Input};
  \node[decision, below of=input] (check) {Valid?};
  \node[block, right of=check, xshift=3cm] (process) {Process};
  \node[block, below of=check, yshift=-1.5cm] (error) {Error};
  \node[block, below of=process] (end) {End};

  \draw[arrow] (start) -- (input);
  \draw[arrow] (input) -- (check);
  \draw[arrow] (check) -- node[anchor=south] {Yes} (process);
  \draw[arrow] (check) -- node[anchor=east] {No} (error);
  \draw[arrow] (process) -- (end);
\end{tikzpicture}
\end{latin}
\caption{فلوچارت ساده}
\label{fig:flowchart}
\end{figure}

%========================================================================
% CHAPTER 8 — CONCLUSION
%========================================================================
\chapter{جمع‌بندی}
\label{chap:conclusion}

این کتاب نقطه‌ی شروعی برای یادگیری برنامه‌نویسی و
الگوریتم بود. مفاهیمی که آموختید:

\begin{itemize}
    \item مبانی پایتون (\cref{chap:basics}).
    \item الگوریتم‌های جستجو (\cref{chap:algorithms}).
    \item ساختمان داده‌های پایه (\cref{chap:datastructures}).
    \item توابع و ماژول‌ها (\cref{chap:functions}).
\end{itemize}

برای ادامه‌ی مسیر، منابع زیر توصیه می‌شوند:
\cite{knuth1984,cormen2009}.

\begin{flushleft}
    \textbf{پایان} \\
    مهر ۱۴۰۵
\end{flushleft}

%------------------------------------------------------------------------
% BACK MATTER (symmetric margins)
%------------------------------------------------------------------------
\backmatter
\frontmattergeometry

% Bibliography — NOT wrapped in \begin{latin}
\printbibliography[title={منابع و مراجع}]

% Index
\printindex

%------------------------------------------------------------------------
% BACK COVER (must be at the very end)
%------------------------------------------------------------------------
\authorbio{%
    تیم توسعه‌ی \lr{MatinBook}، گروهی از پژوهشگران و
    برنامه‌نویسان فارسی‌زبان است که با هدف ارائه‌ی یک
    راه‌حل حرفه‌ای برای حروف‌چینی کتاب‌های فنی فارسی
    فعالیت می‌کند. این پروژه حاصل سال‌ها تجربه در
    حوزه‌ی نشر دیجیتال و برنامه‌نویسی است.
}

\makebackcover{%
    \textbf{هنر برنامه‌نویسی} کتابی جامع برای یادگیری
    اصول برنامه‌نویسی و الگوریتم است.

    \vspace{0.4cm}

    \textbf{در این کتاب می‌آموزید:}
    \begin{itemize}
        \item مفاهیم پایه‌ی برنامه‌نویسی با پایتون
        \item الگوریتم‌های جستجو و مرتب‌سازی
        \item ساختمان داده‌های پرکاربرد
        \item تحلیل پیچیدگی و بهینه‌سازی
        \item برنامه‌نویسی شیءگرا
    \end{itemize}

    \vspace{0.4cm}

    \textbf{ویژگی‌های کتاب:}
    \begin{itemize}
        \item بیش از \pnum{۵۰} مثال کاربردی
        \item تمرین‌های متنوع با پاسخ تشریحی
        \item نمودارها و شبه‌کدهای استاندارد
        \item نمایه‌ی جامع واژگان
    \end{itemize}
}

\end{document}